\section{Introduction}
\label{sec:intro}

Benchmarks saturate quickly. A new test set is released, frontier models climb toward its ceiling within a few years, and soon the top of the leaderboard is a cluster of models separated by fractions of a point~\citep{kiela2021dynabench,ott2022mapping}. Those small gaps still drive model selection, paper acceptances, and research attention. {\bf When do they stop meaning anything?} If we could predict this point from the design of a benchmark, builders could make benchmarks that last longer, and readers would know when to stop trusting a rank difference.

A natural candidate is the {\em item-difficulty distribution}. Item response theory (IRT) tells us that an item discriminates best between models whose ability is close to its difficulty~\citep{lord1980applications,polo2024tinybenchmarks}. A benchmark whose items all sit at one difficulty level should therefore stop separating models once they pass that level, while a benchmark with a long tail of hard items should keep a reserve of informative items for stronger models. This intuition leads to a concrete hypothesis that we test in this paper:
\begin{itemize}[leftmargin=*,itemsep=0pt,topsep=2pt]
    \item {\bf H1 (shape).} Benchmarks with uniform difficulty distributions lose discriminative power once cohort mean accuracy exceeds about 85\%; heavy-tailed distributions keep discriminating up to about 95\%.
    \item {\bf H2 (noise).} At saturation, evaluation-seed variation accounts for more ranking variance than true capability differences.
    \item {\bf H3 (lifetime).} A benchmark's effective lifetime correlates strongly ($r > 0.8$) with the variance of its item difficulties.
\end{itemize}

Prior work has established the pieces but not the link between them. IRT analyses show that LLM benchmarks have strong latent structure and that few items carry most of the information~\citep{polo2024tinybenchmarks,kipnis2024metabench,rodriguez2021evaluation}. Separate work shows that evaluation scores carry non-trivial seed and prompt variance~\citep{madaan2024quantifying,miller2024adding,zheng2024large}. What is missing is a quantitative account of {\em which} benchmark property decides the accuracy at which discrimination collapses, one that separates the shape of the difficulty distribution from its variance and from the number of items, together with a measurement of how seed noise compares with capability differences as a benchmark saturates.

We answer this with four complementary experiments. First, we build a controlled IRT model in which we vary the shape of the difficulty distribution (seven shapes) while holding its variance fixed, and we compute exactly when the ranking of a cohort of similar models becomes dominated by noise. Second, we test the same predictions on real per-item responses of 220 \leaderboard models on 62 tasks, including benchmarks we build by resampling real \mmlu and \hellaswag items into prescribed difficulty shapes. Third, we run 15 open LLMs with 10 option-order seeds on 1{,}500 items to measure seed noise directly. Fourth, we track the real frontier cohort of the leaderboard month by month.

The evidence goes against two parts of the hypothesis and gives only conditional support to the third. With variance matched, all seven difficulty shapes lose discrimination within about 1.5 points of each other (at 95.2--96.6\% accuracy for a 500-item benchmark), and the hard-tailed shape fails slightly {\em earlier}. Real items resampled into the same shapes reproduce this null. What moves the saturation point is the {\em number of items} (80--84\% at 100 items, 96\% at 500, 99\% at 2000) and the {\em spread of the competing models} (83--88\% for a tight cohort, 99\% for a wide one). Difficulty variance predicts lifetime only at fixed size ($r = 0.92$ at $N \ge 500$); the correlation drops to $r = 0.20$ when sizes vary, turns negative at $N = 100$, and is absent across real tasks ($r = -0.06$). Seed noise is about $0.40\times$ the binomial item-sampling noise, explains 1--18\% of score variance on 500-item benchmarks, and dominates only for near-tied models on small subsets at $\ge 99\%$ accuracy. A short calculation (\secref{sec:mechanism}) explains all three findings: at fixed accuracy, the error mass $1 - A$ caps the number of informative items, and shape only decides how much of that capped budget is wasted.

Our contributions are:
\begin{itemize}[leftmargin=*,itemsep=0pt,topsep=0pt]
    \item We separate the effects of difficulty shape, difficulty variance, benchmark size, and cohort spread on the accuracy at which a benchmark stops discriminating, and show that size and cohort spread dominate while shape changes the saturation point by about 1.5 points.
    \item We test the hypotheses on real data from 220 leaderboard models and 62 tasks, including a direct test with real items resampled into prescribed difficulty shapes, and find no relation between difficulty variance and lifetime across real tasks.
    \item We measure evaluation-seed noise for 15 open LLMs with 10 option-order seeds each, and map when it overtakes capability differences as a function of accuracy and item count.
    \item We give a simple approximation that explains these results and translate it into concrete guidance: budget items near the frontier, keep $N$ large, and report seed-level error bars for near-tied models on small, near-ceiling test sets.
\end{itemize}

\Secref{sec:related} reviews related work. \Secref{sec:method} defines our measures and experiments. \Secref{sec:results} reports results for each hypothesis, and \secref{sec:discussion} discusses the mechanism, implications, and limitations.
